\documentclass[aspectratio=169,10pt]{beamer}

% ============================================================================
%  Optimal Transport — Chapter 1 Deck
%  Book: C. Villani, "Optimal Transport: Old and New"
%  Chapter 1: Couplings and Changes of Variables
% ============================================================================

\usepackage[academic]{template-lib/template-lib}
\uselayout{text}
\uselayout{eq}
\uselayout{twocol}
\uselayout{1img}

\TLsetfoot{最优传输 · 第一章：耦合与变量替换}

\title{最优传输：旧与新}
\subtitle{第一章 \enspace 耦合与变量替换}
\author{Cédric Villani}
\institute{Grundlehren der mathematischen Wissenschaften 338\\Springer-Verlag, 2009}
\date{读书报告}

\begin{document}

% ── Title Slide ─────────────────────────────────────────────────────────────
\TLtitlecenter
  {最优传输：旧与新}
  {第一章\enspace 耦合与变量替换}
  {Cédric Villani}
  {读书报告\enspace·\enspace 2026}

% ════════════════════════════════════════════════════════════════════════════
\TLsection{Part I\enspace 耦合 (Couplings)}
% ════════════════════════════════════════════════════════════════════════════

% ── S0: 测度论速成 ────────────────────────────────────────────────────────
\begin{frame}{测度论速成——四个你需要知道的概念}
  \begin{TLtext}
    \begin{description}
      \item[\TLterm{概率测度} $\mu$]
        概率分布的严格形式。你熟悉的 PDF $f(x)$ 定义了
        $\mu(A)=\int_A f(x)\,dx$；PMF 则 $\mu(\{x_i\})=p_i$。
        测度论允许没有密度函数的分布（如 $\delta_x$）。

      \item[\TLterm{可测集} $A$]
        可以定义 $\mu(A)$ 的集合——直觉：``行为良好''的集合。
        你熟悉的区间、球、多面体都是可测的。
        不必深究：所有实际遇到的集合都可测。

      \item[\TLterm{a.s.}\;/\;\TLterm{a.e.}]
        ``几乎必然'' (almost surely) / ``几乎处处'' (almost everywhere)：
        除概率为0的集合外成立。
        类比：连续随机变量取某个特定值的概率为0。

      \item[\TLterm{推前} $T_{\#}\mu$]
        若 $X\sim\mu$，则 $T(X)$ 的分布就是 $T_{\#}\mu$。
        等价：$T_{\#}\mu(B)=\mu(T^{-1}(B))$。
        类比：$Y=g(X)$ 的 PDF 是 $f_Y(y)=f_X(g^{-1}(y))\cdot|g'(g^{-1}(y))|^{-1}$。
    \end{description}
  \end{TLtext}
\end{frame}

% ── S1: 什么是耦合？ ───────────────────────────────────────────────────────
\begin{frame}{耦合的定义 (Definition 1.1)}
  \begin{TLtext}
    设 $(\mathcal{X},\mu)$ 和 $(\mathcal{Y},\nu)$ 为两个概率空间
    （即 $\mu,\nu$ 是两个概率分布）。
    \TLterm{耦合} $\mu$ 与 $\nu$ 意味着构造两个随机变量 $X$ 和 $Y$，
    使得 $\mathrm{law}(X)=\mu$，$\mathrm{law}(Y)=\nu$。
  \end{TLtext}

  \vspace{0.3em}
  等价表述：构造 $\mathcal{X}\times\mathcal{Y}$ 上的联合分布 $\pi$，
  满足\TLterm{边际条件}——联合分布的行和等于 $\mu$，列和等于 $\nu$：
  \[
    \pi[A\times\mathcal{Y}]=\mu[A], \quad
    \pi[\mathcal{X}\times B]=\nu[B]
  \]

  \vspace{0.2em}
  {\small\textbf{2×2 例子：}$X,Y\in\{0,1\}$，$\mu=(\tfrac{1}{2},\tfrac{1}{2})$，
  $\nu=(\tfrac{1}{4},\tfrac{3}{4})$。联合分布\;
  $\pi=\begin{pmatrix}0.10 & 0.40 \\ 0.15 & 0.35\end{pmatrix}$\;
  行和 $=(0.5,0.5)=\mu$，列和 $=(0.25,0.75)=\nu$。\ding{51}}
\end{frame}

% ── S2: 耦合的存在性 ───────────────────────────────────────────────────────
\begin{frame}{耦合的存在性}
  \begin{TLtext}
    \begin{itemize}
      \item \textbf{独立耦合}（最简单）总是存在：令 $X$, $Y$ 独立，
        则联合分布为 $\mu\otimes\nu$\;（即
        $\pi(x,y)=\mu(x)\cdot\nu(y)$）
      \item 另一极端：$Y$ 完全由 $X$ 决定 \enspace $\Longrightarrow$ \enspace
        \TLterm{确定性耦合}
    \end{itemize}
  \end{TLtext}

  \vspace{0.3em}
  {\small\textbf{回到 2×2 例子：}$\mu=(\tfrac{1}{2},\tfrac{1}{2})$，
  $\nu=(\tfrac{1}{4},\tfrac{3}{4})$。\;
  独立耦合\;
  $\pi_{\mathrm{ind}}=\begin{pmatrix}1/8 & 3/8 \\ 1/8 & 3/8\end{pmatrix}$；\;
  确定性耦合 $Y=X\cdot T$\;
  $\pi_{\mathrm{det}}=\begin{pmatrix}0 & 1/2 \\ 1/4 & 1/4\end{pmatrix}$。}

  \vspace{0.3em}
  \TLtakeaway{%
    耦合总是存在的（至少有独立耦合），但确定性耦合不一定存在。%
  }
\end{frame}

% ── S3: 确定性耦合 ─────────────────────────────────────────────────────────
\begin{frame}{确定性耦合 (Definition 1.2)}
  \begin{TLtext}
    耦合 $(X,Y)$ 称为\TLterm{确定性}的，若存在可测函数
    $T\colon\mathcal{X}\to\mathcal{Y}$ 使得 $Y=T(X)$。
    \enspace（直觉：知道了 $X$ 就完全确定了 $Y$。）
  \end{TLtext}

  \vspace{0.3em}
  等价于以下任一表述：
  \begin{itemize}
    \item 联合分布 $\pi$ 集中在 $T$ 的\TLterm{图} (graph) 上
      \;（即 $\pi$ 只在 $\{(x,T(x))\}$ 处有质量）
    \item $T_{\#}\mu = \nu$\enspace （推前：$T$ 把 $\mu$ ``推成''了 $\nu$）
    \item \TLterm{变量替换公式}：$\displaystyle
      \int_{\mathcal{Y}}\varphi(y)\,d\nu(y)
      =\int_{\mathcal{X}}\varphi(T(x))\,d\mu(x)$\;
      （换元积分的抽象形式）
    \item $\pi = (\mathrm{Id}, T)_{\#}\mu$
  \end{itemize}

  \vspace{0.3em}
  \TLalertblock[注意]{%
    确定性耦合不一定存在（如 $\mu=\delta_x$ 而 $\nu$ 非原子），也可能有无穷多个。%
  }
\end{frame}

% ════════════════════════════════════════════════════════════════════════════
\TLsection{Part II\enspace 经典耦合}
% ════════════════════════════════════════════════════════════════════════════

% ── S4: 经典耦合概览 ───────────────────────────────────────────────────────
\begin{frame}{八种经典耦合}
  \begin{TLtext}
    \begin{enumerate}
      \item \TLterm{可测同构} — 无原子 Polish 空间之间的双射
      \item \TLterm{Moser 映射} — 通过椭圆方程构造的光滑传输
      \item \TLterm{单调重排} (Increasing Rearrangement on $\mathbb{R}$)
      \item \TLterm{Knothe–Rosenblatt 重排} (多变量递推)
      \item \TLterm{Holley 耦合} — 格上的随机序耦合
      \item \TLterm{抛物方程的概率表示} — PDE 与 SDE 的桥梁
      \item \TLterm{随机过程的精确耦合} — Markov 链耦合
      \item \TLterm{最优耦合} (Optimal Transport) — Monge–Kantorovich 问题
    \end{enumerate}
  \end{TLtext}
\end{frame}

% ── S5: 可测同构与 Moser 映射 ──────────────────────────────────────────────
\begin{frame}{可测同构与 Moser 映射}
  \begin{TLtext}
    \textbf{可测同构}\enspace 设 $(\mathcal{X},\mu)$, $(\mathcal{Y},\nu)$ 为无原子的
    \TLterm{Polish 概率空间}（完备可分度量空间——$\mathbb{R}^n$ 就是典型例子），
    则存在可测双射 $T\colon\mathcal{X}\to\mathcal{Y}$ 使得 $T_{\#}\mu=\nu$。\enspace
    \textcolor{TLneg}{\ding{55}}\; $T$ 通常非常奇异（不光滑），实用价值有限。

    \medskip
    \textbf{Moser 映射}\enspace 在光滑紧 \TLterm{Riemann 流形}
    （直觉：弯曲空间上的``曲面''，局部看起来像 $\mathbb{R}^n$）上，
    设 $f,g$ 为正的 Lipschitz 概率密度，
    可通过求解椭圆方程构造确定性耦合。\enspace
    \textcolor{TLpos}{\ding{51}}\; $T$ 的正则性与 $f,g$ 一致：
    若 $f,g\in C^{k,\alpha}$，则 $T\in C^{k+1,\alpha}$。
  \end{TLtext}
\end{frame}

% ── S6: 单调重排 ───────────────────────────────────────────────────────────
\begin{frame}{单调重排 (Increasing Rearrangement on $\mathbb{R}$)}
  设 $\mu,\nu$ 为 $\mathbb{R}$ 上的概率测度，定义\TLterm{累积分布函数}：
  \[
    F(x)=\int_{-\infty}^{x}d\mu,\qquad
    G(y)=\int_{-\infty}^{y}d\nu
  \]

  及其\TLterm{右连续逆}：
  \[
    F^{-1}(t)=\inf\{x\in\mathbb{R} : F(x)>t\},\qquad
    G^{-1}(t)=\inf\{y\in\mathbb{R} : G(y)>t\}
  \]

  \TLtakeaway{%
    传输映射 $T=G^{-1}\circ F$ 将 $\mu$ 推前为 $\nu$。
    单调、显式、具有良好的几何性质。%
  }
\end{frame}

% ── S7: Knothe–Rosenblatt ─────────────────────────────────────────────────
\begin{frame}{Knothe–Rosenblatt 重排 ($\mathbb{R}^n$)}
  \begin{TLtext}
    设 $\mu\ll\mathrm{Leb}$\;（即 $\mu$ 有密度函数——``绝对连续于 Lebesgue 测度''），
    $\nu$ 为 $\mathbb{R}^n$ 上的概率测度。
    \TLterm{逐变量递推}构造耦合：
  \end{TLtext}

  \vspace{0.3em}
  \begin{enumerate}
    \item \textbf{Step 1:} 取第一边际 $\mu_1(dx_1)$, $\nu_1(dy_1)$，
      用单调重排定义 $y_1=T_1(x_1)$
    \item \textbf{Step 2:} 对\TLterm{条件分布}
      $\mu_2(dx_2|x_1)$, $\nu_2(dy_2|y_1)$ 定义 $y_2=T_2(x_2;x_1)$
      \;（类比：给定 $X_1=x_1$ 后 $X_2$ 的分布）
    \item[] $\vdots$
    \item[$n$.] \textbf{Step $n$:} 得到映射 $y=T(x)$，将 $\mu$ 传输为 $\nu$
  \end{enumerate}

  \vspace{0.3em}
  \medskip
  \textbf{关键性质：}Jacobi 矩阵 $\nabla T$ 为\TLterm{上三角}且对角元为正——适用于几何应用。
  但不具有坐标不变性。
\end{frame}

% ── S8: Holley 耦合 ───────────────────────────────────────────────────────
\begin{frame}{Holley 耦合 (格上序耦合)}
  \begin{TLtext}
    设 $\mu,\nu$ 为有限格 $\Lambda=\{0,1\}^N$ 上的离散概率分布，
    配备自然偏序 ($x\le y$ $\Leftrightarrow$ $x_n\le y_n,\ \forall n$)。
  \end{TLtext}

  \medskip
  \textbf{FKG 条件 (1.2)：}
  若 $\mu[\inf(x,y)]\;\nu[\sup(x,y)] \ge \mu[x]\;\nu[y]$，
  则存在耦合 $(X,Y)$ 使得 $X\le Y$\;（$\mu+\nu$-a.s.）。\enspace
  直觉：$\nu$ 比 $\mu$ 更集中于``大值''。

  \medskip
  {\small 应用：统计力学中的 FKG 不等式、相变理论。}
\end{frame}

% ── S9: 最优传输 ──────────────────────────────────────────────────────────
\begin{frame}{最优耦合 (Optimal Transport)}
  引入\TLterm{代价函数} $c(x,y)$：将单位质量从 $x$ 移到 $y$ 的``功''
  （直觉：搬家的``路费''）。

  \vspace{0.3em}
  \TLeqsingle{%
    \inf_{\pi\in\Pi(\mu,\nu)}
    \int_{\mathcal{X}\times\mathcal{Y}} c(x,y)\,d\pi(x,y)
  }{%
    \small Monge--Kantorovich 问题：在所有联合分布（\TLterm{传输计划}）中
    最小化总代价。%
  }

  \vspace{0.3em}
  {\small\textbf{离散直觉：}2 个仓库 $\to$ 2 个商店，代价矩阵 $C=\bigl[\begin{smallmatrix}1&3\\3&1\end{smallmatrix}\bigr]$。\\
  方案 A（确定性）：$1\!\to\!1,\;2\!\to\!2$，总代价 $=1+1=2$。\;
  方案 B（``混合''）：$1\!\to\!2,\;2\!\to\!1$，总代价 $=3+3=6$。\;
  方案 A 更优。}

  \medskip
  \textbf{Monge 问题：}在一定条件下，最优耦合是\TLterm{确定性}的
  （即存在传输映射 $T$，不分裂质量——每个仓库只送一家）。
  典型结果：$c(x,y)=|x-y|^2$，$\mu\ll\mathrm{Leb}$，
  有限二阶矩 $\Rightarrow$ 存在唯一最优 Monge 耦合。
\end{frame}

% ── S10: 最优耦合的性质 ───────────────────────────────────────────────────
\begin{frame}{最优耦合的性质}
  \begin{TLtext}
    \begin{itemize}
      \item[\TLyes] 自然出现在经济学、物理学、PDE、几何中的许多问题
      \item[\TLyes] 对扰动具有良好的\TLterm{稳定性}
      \item[\TLyes] 编码了底层的几何信息（当 $c$ 由几何定义时）
      \item[\TLyes] 在光滑和非光滑设定中均存在
      \item[\TLyes] 具有丰富的结构：对偶变分问题、连续插值
      \item[\TLno] 一般而言\TLterm{正则性有限}——即使对很好的代价函数也存在反例
    \end{itemize}
  \end{TLtext}

  \vspace{0.3em}
  \TLtakeaway{%
    单调重排和 Holley 耦合都可以视为最优传输的特例。%
  }
\end{frame}

% ════════════════════════════════════════════════════════════════════════════
\TLsection{Part III\enspace 粘合引理与变量替换}
% ════════════════════════════════════════════════════════════════════════════

% ── S11: 粘合引理 ─────────────────────────────────────────────────────────
\begin{frame}{粘合引理 (Gluing Lemma)}
  \begin{TLtext}
    \textbf{直观：}若 $Z$ 是 $Y$ 的函数，$Y$ 是 $X$ 的函数，则 $Z$ 也是 $X$ 的函数。
    这个性质在\TLterm{非确定性}耦合中仍然成立！
  \end{TLtext}

  \vspace{0.5em}
  \TLinfoblock{粘合引理}{%
    设 $(\mathcal{X}_i,\mu_i)$，$i=1,2,3$，为 Polish 概率空间。
    若 $(X_1,X_2)$ 是 $(\mu_1,\mu_2)$ 的耦合，
    $(Y_2,Y_3)$ 是 $(\mu_2,\mu_3)$ 的耦合，则可构造三元组
    $(Z_1,Z_2,Z_3)$ 使得
    \begin{itemize}
      \item $(Z_1,Z_2)$ 与 $(X_1,X_2)$ 同分布
      \item $(Z_2,Z_3)$ 与 $(Y_2,Y_3)$ 同分布
    \end{itemize}%
  }

  \vspace{0.3em}
  {\small\textbf{构造：}将 $\pi_{12}$ 与 $\pi_{23}$ 沿公共边际 $\mu_2$ ``粘合''——就像沿公共边拼两块拼图。}
\end{frame}

% ── S12: 粘合引理的公式 ───────────────────────────────────────────────────
\begin{frame}{粘合引理——分解与重构}
  分解两个联合分布（\TLterm{分解/disintegration}——即``已知 $x_2$ 时 $x_1$ 的条件分布''）：
  \[
    \pi_{12}(dx_1\,dx_2) = \underbrace{\pi_{12}(dx_1|x_2)}_{\text{条件分布}}\;\mu_2(dx_2)
  \]
  \[
    \pi_{23}(dx_2\,dx_3) = \underbrace{\pi_{23}(dx_3|x_2)}_{\text{条件分布}}\;\mu_2(dx_2)
  \]

  重构三元联合分布：
  \[
    \boxed{\pi_{123}(dx_1\,dx_2\,dx_3)
      = \pi_{12}(dx_1|x_2)\;\mu_2(dx_2)\;\pi_{23}(dx_3|x_2)}
  \]

  {\small 直觉：已知 $x_2$ 后，$x_1$ 与 $x_3$ 条件独立，所以可以``粘合''。}

  \TLtakeaway{%
    粘合引理是构建多边际耦合的基本工具，在最优传输理论中反复使用。%
  }
\end{frame}

% ════════════════════════════════════════════════════════════════════════════
\TLsection{Part IV\enspace 变量替换与守恒律}
% ════════════════════════════════════════════════════════════════════════════

% ── S13: 变量替换公式 ─────────────────────────────────────────────────────
\begin{frame}{变量替换公式}
  设 $M$ 为 $n$ 维 Riemann 流形，$\mu_0,\mu_1$ 为概率测度，
  $T\colon M\to M$ 可测且 $T_{\#}\mu_0=\mu_1$。

  \vspace{0.3em}
  \textbf{假设：}
  \begin{enumerate}
    \item $\mu_0(dx)=\rho_0(x)\,\nu(dx)$, \quad
      $\mu_1(dy)=\rho_1(y)\,\nu(dy)$\quad（$\nu=e^{-V}\mathrm{vol}$）
    \item $T$ 是\TLterm{单射}（一一对应，不``折叠''空间）
    \item $T$ 是\TLterm{局部 Lipschitz} 的
      \;（直觉：``不会无限陡峭''——变化率有界）
  \end{enumerate}

  \vspace{0.3em}
  则 $\mu_0$-a.s.（即``除了概率为零的集合外处处成立''）：
  \[
    \boxed{\rho_0(x) = \rho_1(T(x))\;J_T(x)}
  \]

  其中 $J_T(x) = \displaystyle\lim_{\varepsilon\downarrow 0}
  \frac{\nu[T(B_\varepsilon(x))]}{\nu[B_\varepsilon(x)]}$ 为\TLterm{Jacobi 行列式}
  \;（直觉：$T$ 把 $x$ 处的小球体积放大/缩小了多少倍）。
\end{frame}

% ── S14: 变量替换——推广 ──────────────────────────────────────────────────
\begin{frame}{变量替换公式的推广}
  \textbf{正则性要求}\enspace 条件 (iii) 可放宽为：
  $T$ 是\TLterm{近似可微}的 (approximately differentiable)。
  \;（直觉：``几乎处处''可微——不可微点构成一个``可忽略''的集合。）

  \medskip
  \textbf{实用结论}\enspace 变量替换公式在以下情况成立：
  \begin{itemize}
    \item[\textcolor{TLpos}{\ding{51}}] $T$ 为 Lipschitz 映射
    \item[\textcolor{TLpos}{\ding{51}}] $T$ 为\TLterm{Sobolev 函数}
      \;（允许``几乎处处''可微的广义可微函数）
    \item[\textcolor{TLpos}{\ding{51}}] \TLterm{BV 函数}（有界变差函数）
      \;（导数为测度，允许跳跃间断）
    \item[\textcolor{TLpos}{\ding{51}}] $T=\nabla\varphi$（凸函数梯度）
      \;（凸函数天然几乎处处可微）
  \end{itemize}
\end{frame}

% ── S14.5: 微分算子速成 ──────────────────────────────────────────────────
\begin{frame}{微分算子速成}
  你已知梯度 $\nabla f$（方向导数），下面两个是它的``亲戚''：

  \vspace{0.5em}
  \textbf{1. 散度 (divergence)} $\nabla\cdot\xi$
  \[
    \nabla\cdot\xi = \frac{\partial\xi_1}{\partial x_1}
      + \frac{\partial\xi_2}{\partial x_2} + \cdots
      + \frac{\partial\xi_n}{\partial x_n}
  \]
  {\small 直觉：速度场在某点是``源''（$>0$）还是``汇''（$<0$）。\quad
    类比：水流——某点不断有水涌出 = 正散度。}

  \vspace{0.8em}
  \textbf{2. Laplacian} $\Delta f = \nabla\cdot(\nabla f)$
  \[
    \Delta f = \frac{\partial^2 f}{\partial x_1^2}
      + \frac{\partial^2 f}{\partial x_2^2} + \cdots
      + \frac{\partial^2 f}{\partial x_n^2}
  \]
  {\small 直觉：$\Delta f$ 衡量 $f$ 在某点比周围平均``凸''还是``凹''。\quad
    例：$\Delta f>0$ $\Rightarrow$ 该点是局部极小值附近。}
\end{frame}

% ── S15: 质量守恒 ─────────────────────────────────────────────────────────
\begin{frame}{质量守恒公式}
  连续物理中最重要的变量替换定理源自\TLterm{质量守恒方程}：

  \vspace{0.5em}
  \TLeqsingle{%
    \frac{\partial\rho}{\partial t} + \nabla\cdot(\rho\,\xi) = 0
  \quad\text{(1.5)}}
  {\small
    $\rho(t,x)$：粒子密度 \quad $\xi(t,x)$：速度场 \quad
    $\nabla\cdot$：散度（``源''还是``汇''）}

  \vspace{0.5em}
  \textbf{Euler $\leftrightarrow$ Lagrange 类比：}
  \begin{itemize}
    \item \textbf{Euler 描述}（场）：站在河边看水流——观测密度/速度场
    \item \textbf{Lagrange 描述}（粒子）：坐在漂流船上——追踪单个粒子的轨迹
  \end{itemize}
  {\small 两者等价：知道所有粒子的轨迹 $\Leftrightarrow$ 知道每一时刻的密度场。}
\end{frame}

% ── S16: 质量守恒定理 ─────────────────────────────────────────────────────
\begin{frame}{质量守恒定理——精确表述}
  \TLinfoblock{定理}{%
    设 $M$ 为 $C^1$ 流形，$\xi(t,x)$ 为可测速度场，
    $(\mu_t)_{0\le t<T}$ 为弱连续的概率测度族，满足
    $\displaystyle\int_0^T\!\int_M|\xi(t,x)|\,\mu_t(dx)\,dt<+\infty$。
    则以下等价：\\[0.5em]
    \textbf{(i)} $\mu_t$ 是线性传输 PDE
      $\partial_t\mu+\nabla_x\cdot(\mu\,\xi)=0$ 的\TLterm{弱解}
      \;（``用测试函数验证''的广义解，不要求可微）；\\[0.3em]
    \textbf{(ii)} $\mu_t$ 是随机解 $T_t(x)$ 在时刻 $t$ 的分布。\\[0.5em]
    若 $\xi$ 局部 Lipschitz，则 (ii) 强化为：$\mu_t=(T_t)_{\#}\mu_0$。%
  }

  \vspace{0.3em}
  \TLtakeaway{%
    即使 $\xi$ 无正则性，质量守恒仍成立——此时流可能不唯一，需考虑随机解。%
  }
\end{frame}

% ── S17: 扩散公式 ─────────────────────────────────────────────────────────
\begin{frame}{扩散公式 (Diffusion Theorem)}
  \begin{TLtext}
    与 \TLterm{Itô 公式} 相关（随机微积分的``链式法则''），
    由 Bachelier、Einstein、Smoluchowski 独立发现。
  \end{TLtext}

  \vspace{0.3em}
  设 $M$ 为 \TLterm{Ricci 曲率}下方有界的 Riemann 流形
  \;（直觉：曲率描述空间``弯曲''程度，下界保证粒子不会``跑丢''），
  $X_t$ 满足\TLterm{随机微分方程 (SDE)}：
  \[
    dX_t = \sqrt{2}\;\sigma(t,X_t)\,dB_t
    \qquad (0\le t<T)
  \]
  {\small $dB_t$：Brown 运动的随机增量 \quad
    $\sigma$：扩散系数矩阵 \quad
    $\sigma\,\sigma^*$：\TLterm{扩散张量}（对称正半定矩阵）}

  \vspace{0.3em}
  则以下等价：
  \begin{itemize}
    \item[\textbf{(i)}] $\mu_t$ 是线性扩散 PDE 的弱解：
      $\displaystyle\partial_t\mu = \nabla_x\cdot\bigl((\sigma\,\sigma^*)\nabla_x\mu\bigr)$
    \item[\textbf{(ii)}] $\mu_t = \mathrm{law}(X_t)$
      \;（$X_t$ 的分布）
  \end{itemize}
\end{frame}

% ── S18: 扩散公式——例子 ──────────────────────────────────────────────────
\begin{frame}{扩散公式——例子与注记}
  \textbf{例 1.5}\enspace 在 $\mathbb{R}^n$ 中，\TLterm{热方程}
  $\partial_t\rho=\Delta\rho$
  \;（$\Delta$ = Laplacian，见 S14.5）
  以 $\delta_0$（集中在原点的概率分布）为初值的解是
  $X_t=\sqrt{2}\,B_t$ 的分布（加速 $\sqrt{2}$ 倍的 Brown 运动）。

  \medskip
  \textbf{注记 1.6}\enspace 更精细的判据：只要 $\mathrm{Ric}_x\ge -C\,d(x_0,x)^2\,g_x$（$x\to\infty$），
  扩散方程就成立。指数 $2$ 是最优的。

  \vspace{0.3em}
  \TLalertblock[练习 1.7]{%
    在紧流形上，利用热方程解 $(\rho_t)_{t\ge 0}$ 构造 $\rho_0$ 与 $\rho_1$
    的确定性耦合。提示：将热方程改写为质量守恒方程。%
  }
\end{frame}

% ════════════════════════════════════════════════════════════════════════════
\TLsection{Part V\enspace Moser 耦合 (附录)}
% ════════════════════════════════════════════════════════════════════════════

% ── S19: Moser 耦合构造 ──────────────────────────────────────────────────
\begin{frame}{Moser 耦合——构造}
  设 $M$ 为光滑 $n$ 维 Riemann 流形，$\nu(dx)=e^{-V(x)}\mathrm{vol}(dx)$，
  $\mu_0=\rho_0\,\nu$, $\mu_1=\rho_1\,\nu$。

  \vspace{0.3em}
  \textbf{步骤：}
  \begin{enumerate}
    \item 求解\TLterm{椭圆方程}：
      $\displaystyle(\Delta-\nabla V\cdot\nabla)\,u = \rho_0-\rho_1$
      \;（直觉：``稳态''方程——不含时间导数，描述平衡状态）
    \item 定义速度场：
      $\displaystyle\xi(t,x)=\frac{\nabla u(x)}{(1-t)\,\rho_0(x)+t\,\rho_1(x)}$
    \item 令 $\mu_t=(1-t)\mu_0+t\mu_1$，则：
  \end{enumerate}

  \[
    \partial_t\mu = (\rho_1-\rho_0)\,\nu,\qquad
    \nabla\cdot(\mu_t\,\xi)=(\rho_0-\rho_1)\,\nu
  \]

  \vspace{0.2em}
  由质量守恒 $\Rightarrow$ $\mu_t=(T_t)_{\#}\mu_0$，特别地 $T_1$ 将 $\mu_0$ 推前为 $\mu_1$。
\end{frame}

% ── S20: Moser 耦合——正则性 ──────────────────────────────────────────────
\begin{frame}{Moser 耦合——正则性与适用范围}
  \textbf{紧流形 ($V=0$)}\enspace 若 $\rho_0,\rho_1$ 为正的 Lipschitz 密度：
  \begin{itemize}
    \item $\Delta u = \rho_0-\rho_1$ 的解 $u\in C^{2,\alpha}$
    \item $\Rightarrow \nabla u\in C^{1,\alpha}$
    \item $\Rightarrow T$ 具有良好正则性
  \end{itemize}

  \medskip
  \textbf{一般情形}\enspace 正则性取决于：
  \begin{itemize}
    \item $V$ 的正则性
    \item $V$ 在无穷远处的行为
    \item $\rho_0,\rho_1$ 的下界
  \end{itemize}

  \vspace{0.5em}
  \TLtakeaway{%
    Moser 耦合\textbf{简单、优雅、强大}，在几何中扮演重要角色。
    但需要密度为正且满足正则性条件。%
  }
\end{frame}

% ════════════════════════════════════════════════════════════════════════════
\TLsection{总结}
% ════════════════════════════════════════════════════════════════════════════

% ── S21: 第一章总结 ───────────────────────────────────────────────────────
\begin{frame}{第一章总结}
  \begin{TLtext}
    \begin{itemize}
      \item \TLterm{耦合}是概率论的基本工具——联合分布的构造
      \item 八种经典耦合各有适用场景：
        \begin{itemize}
          \item 单调重排：$\mathbb{R}$ 上简单显式
          \item Knothe–Rosenblatt：$\mathbb{R}^n$ 上的递推构造
          \item \textbf{最优传输}：几何意义丰富，本书的核心
        \end{itemize}
      \item 三大分析工具：
        \begin{itemize}
          \item \TLterm{粘合引理}——多边际耦合的构建
          \item \TLterm{变量替换公式}——Jacobi 行列式
          \item \TLterm{质量守恒}与\TLterm{扩散公式}——PDE 与 SDE 的桥梁
        \end{itemize}
      \item Moser 耦合——光滑情形下的优雅构造
    \end{itemize}
  \end{TLtext}

  \vspace{0.3em}
  \TLtakeaway{%
    第一章为后续的最优传输理论奠定了概率论与分析工具的基础。%
  }
\end{frame}

% ════════════════════════════════════════════════════════════════════════════
\TLsection{附录\enspace 术语速查}
% ════════════════════════════════════════════════════════════════════════════

% ── S22: 术语速查 (上) ────────────────────────────────────────────────────
\begin{frame}{术语速查——概率与分析}
  \begin{columns}[T]
    \begin{column}{0.48\textwidth}
      \textbf{概率论}
      \begin{itemize}\small
        \item \textbf{耦合 (coupling)}：两个分布的``联合分布''
        \item \textbf{推前 (push-forward)}：$T_{\#}\mu$ = 把 $\mu$ 用 $T$ ``推过去''
        \item \textbf{测度 (measure)}：``广义体积/概率''
        \item \textbf{绝对连续 ($\mu\ll\nu$)}：$\mu$ 可由密度 $\rho$ 表示
        \item \textbf{a.s.}：almost surely，``概率为1''
        \item \textbf{弱解}：用测试函数验证的广义解
      \end{itemize}
    \end{column}
    \begin{column}{0.48\textwidth}
      \textbf{分析}
      \begin{itemize}\small
        \item \textbf{Lipschitz}：变化率有界，``不会无限陡峭''
        \item \textbf{Jacobi 行列式}：映射对局部体积的放大率
        \item \textbf{散度 $\nabla\cdot$}：速度场的``源/汇''强度
        \item \textbf{Laplacian $\Delta$}：$\nabla\cdot\nabla$，描述``凸凹''
        \item \textbf{Sobolev}：允许几乎处处可微的广义函数
        \item \textbf{BV}：有界变差，允许跳跃间断
      \end{itemize}
    \end{column}
  \end{columns}
\end{frame}

% ── S23: 术语速查 (下) ────────────────────────────────────────────────────
\begin{frame}{术语速查——几何与方程}
  \begin{columns}[T]
    \begin{column}{0.48\textwidth}
      \textbf{几何}
      \begin{itemize}\small
        \item \textbf{Riemann 流形}：带度量（可测距离/角度）的光滑空间
        \item \textbf{Ricci 曲率}：描述空间弯曲程度（下界保证粒子不跑丢）
        \item \textbf{Polish 空间}：完备可分的度量空间
        \item \textbf{紧流形}：有限大小、无边界的流形
      \end{itemize}
    \end{column}
    \begin{column}{0.48\textwidth}
      \textbf{方程类型}
      \begin{itemize}\small
        \item \textbf{椭圆方程}：稳态方程（无时间导数），如 $\Delta u = f$
        \item \textbf{抛物方程}：扩散/热方程，如 $\partial_t\rho=\Delta\rho$
        \item \textbf{双曲方程}：波动/传输方程，如 $\partial_t\rho+\nabla\cdot(\rho\xi)=0$
        \item \textbf{SDE}：随机微分方程，含随机增量 $dB_t$
      \end{itemize}
    \end{column}
  \end{columns}
\end{frame}

\end{document}
